\documentclass[leqno, a4paper, 12pt]{jsarticle}
\usepackage{amssymb}
\usepackage{amsthm}
\usepackage{amsmath}
\usepackage{comment}
\usepackage{url}

\usepackage{enumitem}
%\usepackage{showkeys}


%定理環境 thmはあまり使わずpropをなるべく使う。
%主張は基本的に証明の中で使い、そのため番号を付けない。
%注意環境を付けてみた。
\newtheorem{thm}{定理}
\newtheorem{prop}[thm]{命題}
\newtheorem{cor}[thm]{系}
\newtheorem{lem}[thm]{補題}
\newtheorem{df}[thm]{定義}
\newtheorem{exam}[thm]{例}
\newtheorem{fact}[thm]{事実}
\newtheorem*{claim}{主張}
\newtheorem*{cau}{注意}
\newtheorem*{rmk}{注意}
\renewcommand{\proofname}{証明}
%矢印たち。
%\toは\rightarrowと同じ。\getsは\leftarrowと同じ。同値は\iff
\newcommand{\To}{\Longrightarrow}
\newcommand{\Gets}{\LongLeftarrow}
%黒板太字
\newcommand{\nn}{\mathbb{N}}
\newcommand{\zz}{\mathbb{Z}}
\newcommand{\qq}{\mathbb{Q}}
\newcommand{\rr}{\mathbb{R}}
\newcommand{\cc}{\mathbb{C}}
%無限大
\newcommand{\mgn}{\infty}
%差集合は\setminusで統一する。等号の否定は\neq
%真部分集合は\subsetneqq　偏微分は\partial
%ディスプレイスタイル。\dfracでディスプレイ分数
\newcommand{\dpst}{\displaystyle}

\title{論文メモ
\large{\\--- その他もろもろ：
層の話とか位相群の帰納極限の話 ---}
}
\author{ナンブ　キトラ}
\date{\today}
\begin{document}
\maketitle

本棚を整理したいので，
溜まっている論文のメモを書いて未来への手紙とすることにする．



論文
\cite{wa1}
は
一般位相空間論の研究者である
児玉之宏による「層の理論」
に関する解説である．
おそらく一般位相空間論と
(現代的な？)代数的トポロジーを結ぶ
歴史的重要文献のような気がするが
あんまりよくわかっていない．




論文
\cite{wa2}
と
\cite{wa3}
は位相群の帰納的極限を考えると
群にはなるのだが，
その帰納極限としての
位相は位相群になるとは限らないという話にまつわるものである．
論文\cite{wa2}によればなんかミルナーの
Introduction to algebraic K-theory 
で$GL_{n}(\cc)$の
帰納的極限を考える議論があるらしい
(一応回避可能らしい)．
しかし一般線形群ぐらいだと
ちゃんと帰納極限が位相群になるともいう．
論文\cite{wa3}に色々書いてあるらしい．


\begin{thebibliography}{99}
\bibitem{wa1}
児玉之宏, 
「Sheaf理論の位相空間への応用について」, 
数理解析研究所講究録, 
148 (1972), 
8--19, 
http://hdl.handle.net/2433/106767

\bibitem{wa2}
山崎愛一，
「一般線形群の帰納極限について」, 
数理解析研究所講究録, 
1017 
(1997), 
92--103, 
http://hdl.handle.net/2433/61637

\bibitem{wa3}
辰馬伸彦
「位相群の帰納的極限の群位相」, 
雑誌「数学」, 
50 
(4)
(1998), 
428--431, 
DOI: 
10.11429/sugaku1947.50.428

\bibitem{wa4}
平井 武 , 
下村 宏彰, 
 辰馬 伸彦, 
 「位相空間の帰納的極限の位相」，
表現論シンポジウム講演集, 
1998, 
63-86, 
DOI：
\verb|10.34508/repsympo.1998.0_63|

\end{thebibliography}



\end{document}